% !TEX root = ../../hw07.tex
% Source: author's supplied text in chat, 2026-11-06, Q7 (i)--(iii).
\begin{proof}[Solution]
  All three quotients are of $0/0$ form as
  $x\to0$. We apply L'H\^opital's rule on both sides of zero.

  For part (i),
  \[
    \begin{aligned}
      \lim_{x\to0}\frac{\sin(ax)}x
        &\overset{\mathrm{L'H}}{=}
          \lim_{x\to0}\frac{a\cos(ax)}1\\
        &=a\cos(0)\\
        &=a.
    \end{aligned}
  \]

  For part (ii), take $a\ne0$.
  On a sufficiently small neighborhood of
  zero, $1+ax>0$ and the denominator derivative
  $a/(1+ax)$ is nonzero. Thus
  \[
    \begin{aligned}
      \lim_{x\to0}\frac{x}{\ln(1+ax)}
        &\overset{\mathrm{L'H}}{=}
          \lim_{x\to0}\frac1{\frac{a}{1+ax}}\\
        &=\lim_{x\to0}\frac{1+ax}{a}\\
        &=\frac1a.
    \end{aligned}
  \]

  For part (iii), applying L'H\^opital's rule twice gives
  \[
    \begin{aligned}
      \lim_{x\to0}\frac{1-\cos(2x)}{\sin^2x}
        &\overset{\mathrm{L'H}}{=}
          \lim_{x\to0}\frac{2\sin(2x)}{2\sin x\cos x}\\
        &\overset{\mathrm{L'H}}{=}
          \lim_{x\to0}\frac{4\cos(2x)}
            {2(\cos^2x-\sin^2x)}\\
        &=\frac42\\
        &=2.
    \end{aligned}
  \]
  The intermediate quotient is again of $0/0$
  form. The two denominator derivatives are $2\sin x\cos x$ and
  $2(\cos^2x-\sin^2x)$. They are nonzero on a sufficiently small
  punctured neighborhood of zero, and the final derivative ratio
  tends to $2$, justifying both applications.
\end{proof}
